\section{Limitations and Future Work}

The current evaluation has several limitations that scope the conclusions. First, the proof-of-concept experiments use a single operating point (simulation hour 7) with favorable attack conditions; scenarios with lower base load, higher thresholds, or faster controller response would better differentiate attacker capabilities and more clearly reveal the evaluation validity gap. Second, the 5-minute experiment window captures only short-term attack-response dynamics; longer campaigns spanning 1--2 hours would test resource management strategies and reveal whether the efficiency advantages of adaptive attackers compound over time. Third, the threshold-based defender is simpler than production systems, which may incorporate anomaly detection, predictive control, or coordinated protection schemes. Fourth, the current primitive catalog is limited to EV capacity manipulation; real adversaries may combine load-altering attacks with false data injection or communication disruption. Finally, the software co-simulation cannot fully capture the timing constraints and communication latencies of real power system hardware.

Several extensions are planned: (1) replacing scripted policies with fully autonomous LLM agents using multi-agent architectures similar to PentestGPT and AttackLLM \cite{deng2024pentestgpt,attackllm}, with retrieval-augmented generation supplying power system models and attack templates; (2) integrating data-driven and physics-informed intrusion detection systems to quantify $\Delta E(M)$ for detection mechanisms and enable controlled comparisons between RL-based and LLM-based attackers \cite{aoufi2020fdiasurvey,sahani2023smartgridids}; (3) expanding the primitive catalog to include false data injection, malicious SCADA commands, and communication disruption for multi-stage attack evaluation; (4) using LLM-GridEval to generate diverse synthetic attack traces that capture adaptive adversary behaviors for IDS benchmarking, complementing established static datasets; and (5) validating the framework on hardware-in-the-loop (HIL) testbeds such as OPAL-RT to evaluate LLM-based attacks against real-time power system simulators with realistic timing constraints and communication latencies. We also plan to release the framework as open-source software to support reproducible security research in the smart grid community.

\section{Conclusion}

This paper introduced LLM-GridEval, a framework that embeds adaptive, tool-using LLM attackers into HELICS-based power grid co-simulations through a schema-constrained primitive interface. Our proof-of-concept evaluation demonstrated that timing intelligence alone is insufficient for effective attacks; LLM-based adversaries require explicit domain model awareness to outperform random baselines. The V1 to V2 attacker evolution illustrates how the framework enables systematic study of adaptive attack strategies under controlled conditions. These results underscore that static evaluations may underestimate adversarial risk in API-driven, EV/DER-rich grids, motivating further research on LLM-based attackers, advanced defenders, and richer attack scenarios within physically grounded evaluation environments.
